\section{Log-Time Cumulative Distribution Function}
\label{apd:ltcdf}

% Given we have $k$ event matches combined across case reports, we want to compute a measure of the discrepancy between the predicted times
% $t^p_i$ and the manual reference times $t^r_i$ for $i = {1,\dots,k}$ that captures the quality of time recovery across time scales.

% \subsection{Log-time CDF}
Recall the log-time cumulative distribution function is given as follows:
% % \subsection*{1. Truncated Discrepancies}
% For each matched event $i$, compute the discrepancy $s_i = |t^p_i - t^r_i|$. then truncate at the maximum cutoff $S_{\text{max}}$ (\emph{e.g.}, 1 year = 8760 hours):
% \[
%    s'_i = \min(s_i, S_{\text{max}})
% \]
% % \subsection*{2. Log-Transformed Discrepancies}
% Apply the log transformation:
% \[
%     x_i = \log(1 + s'_i)
% \] % so that $s_i=0 \implies x_i=0$.
% % \subsection*{3. Sorted Log-Transformed Discrepancies}
% Sort $x_i$ in ascending order:
% \[
%    x_{(1)} \leq x_{(2)} \leq \cdots \leq x_{(k)}
% \]

% % \subsection*{4. Log-Time Cumulative Distribution Function (CDF)}
% The CDF $F(x)$ is defined as the fraction of log-transformed discrepancies less than or equal to $x$:
\[
   F(x) = \frac{1}{k} \sum_{i=1}^k \mathbf{1}_{\{x_{(i)} \leq x\}}
\]
where $\mathbf{1}$ is the indicator function.

% This results in a step function where:
% \begin{itemize}
%    \item For $x < x_{(1)}$, $F(x) = 0$.
%    \item For $x_{(i)} \leq x < x_{(i+1)}$, $F(x) = \frac{i}{k}$.
%    \item For $x \geq x_{(k)}$, $F(x) = 1$.
% \end{itemize}

% \section*{Properties of the Log-Time Cumulative Distribution}
% \begin{itemize}
%    \item \textbf{Domain}: $x \in [0, \log(1 + S_{\text{max}})]$.
%    \item \textbf{Non-Decreasing}: $F(x)$ increases monotonically with $x$.
%    \item \textbf{Right-Continuous}: Jumps occur at each $x_{(i)}$, with $F(x)$ constant between jumps.
   % \item \textbf{Boundary Values}: 
   % \begin{itemize}
% Then, we get the following boundary values:
% $F(0) \geq 0$ (depending on the smallest discrepancy) and $F(\log(1 + S_{\text{max}})) = 1$.
%    % \end{itemize}
% % \end{itemize}

% % \section*{Area Under the Log-Time Curve (AULTC)}
% With the the CDF defined, 
We compute the AULTC as the area under $F(x)$ from $x = 0$ to $x = \log(1 + S_{\text{max}})$, normalized by $\log(1 + S_{\text{max}})$:
\begin{align*}
\text{AULTC} &= \frac{1}{\log(1 + S_{\text{max}})} \Bigg[ \sum_{i=1}^k (x_{(i)} - x_{(i-1)}) \frac{i}{k} + (\log(1 + S_{\text{max}}) - x_{(k)}) \times 1 \Bigg]
\end{align*}
where $x_{(0)} = 0$. With this definition, AULTC = 1 indicates that discrepancies are zero (perfect recovery), resulting in maximum area $\log(1 + S_{\text{max}})$, and AULTC = 0 indicates that all discrepancies exceed $S_{\text{max}}$, yielding zero area. 
The code for the AULTC calculation is given in the function \texttt{ecdf\textunderscore auc2(\dots)} in the Supplementary Materials.


\begin{remark}
The time unit and cutoff $S_\text{max}$ affect the AULTC calculation, so they must be specified when reporting the AULTC.
\end{remark}
% We note that both the choice of time units and cutoff $S_\text{max}$ affect the AULTC calculation.
Unless otherwise stated, we set $S_{\max}=8760$ hours, corresponding to one year.
This upper limit was chosen because it captures clinically meaningful discrepancies from hours to months while preventing very long-horizon case-history events from dominating the metric.

Because of the $\log(1 + \cdot)$ transformation (rather than the $\log(\cdot)$ transformation which is undefined for zero time-error discrepancy), the discrepancies are non-linearly shifted in the log scale.  The non-linear shift adjusts the relative widths of the step function (particularly for small discrepancies) changing the area calculation.  The cutoff $S_\text{max}$ affects the normalization factor.  Therefore, these values should be chosen based on practicalities for the application.  In our case, we chose hours at the time unit because the only sub-hour descriptions reported were several Apgar scores at the minute-level.

\begin{remark}
    The average log-time discrepancy is non-convex.
\end{remark}
This can be seen from the observation that the $\log(1+\cdot)$ function is ``spikier'' than the L1 function. More formally, the log-time discrepancy has a higher curvature around zero error compared to the L1 loss. Consider the second derivatives of the loss functions with respect to the discrepancies $s = |t^p - t^r|$.  For L1 loss $L_1(s) = s$, the second derivative $\frac{d^2 L_1}{ds^2} = 0$ for $s > 0$.
For the log-time loss, $L_{log}(s) = \log(1 + s)$, the second derivative is $\frac{d^2 L_{log}}{ds^2} = -\frac{1}{(1 + s)^2}$. At $s = 0$, the second derivative is $-1$, while the second derivative of the L1 loss is $0$.

To provide a simple counterexample, define the average log-time discrepancy is given by:
\[
L(t^p) = \frac{1}{k} \sum_{i=1}^k \log(1 + \min(|t^p_i - t^r_i|, S_{\text{max}}))
\]
To prove that this function is non-convex, we need to show that there exist $t^{p1}, t^{p2}$ and $\lambda \in (0,1)$ such that
\[
L(\lambda t^{p1} + (1-\lambda) t^{p2}) > \lambda L(t^{p1}) + (1-\lambda) L(t^{p2})
\]
Consider the case where $k=1$, $t^r = 0$, and $S_{\text{max}} = 2$. The loss function is $L(t^p) = \log(1 + \min(|t^p|, 2))$.
Let $t^{p1} = 0$ and $t^{p2} = 3$. Choose $\lambda = 0.5$.
Then $\lambda t^{p1} + (1-\lambda) t^{p2} = 0.5 \times 0 + 0.5 \times 3 = 1.5$.
Now we evaluate the loss function at these points:
\begin{align*}
L(t^{p1}) &= L(0) = \log(1 + \min(|0|, 2)) = 0 \\
L(t^{p2}) &= L(3) = \log(1 + \min(|3|, 2)) = \\
    & \qquad \log(3) \\
L(\lambda t^{p1} + (1-\lambda) t^{p2}) &= L(1.5) = \log(2.5)
\end{align*}
Now we check the convexity condition:
\[
L(1.5) \leq 0.5 L(0) + 0.5 L(3)
\]
\[
\log(2.5) \leq 0.5 \times 0 + 0.5 \times \log(3)
\]
\[
\log(2.5) \leq \log(\sqrt{3})
\]
Since $2.5 > \sqrt{3} \approx 1.732$, we have $\log(2.5) > \log(\sqrt{3})$.
Therefore,
\[
L(\lambda t^{p1} + (1-\lambda) t^{p2}) > \lambda L(t^{p1}) + (1-\lambda) L(t^{p2})
\]
This proves that the average log-time discrepancy is non-convex.
This remark illustrates the loss form, in case it is considered for optimization/model training rather than for assessment purposes.
% \subsection{Approximation for a shift in units (and lower bound)}

% In this subsection, we provide an approximate proportionate relationship between the \textbf{Area Under the Log-Time Curve (AULTC)} when converting between different units, we introduce the following variables. 

% Let $U$ and $V$ be two units of time with a conversion factor $c$, such that:
% \[
% 1 \text{ unit of } U = c \text{ units of } V
% \]
% (Example: $U = \text{hours}, V = \text{seconds}, c = 3600$).

% Let $S_{\text{max}}^{U}$ be the maximum discrepancy cutoff in units $U$.  Correspondingly, the maximum cutoff in units $V$ is:
% \[
% S_{\text{max}}^{V} = c \times S_{\text{max}}^{U}
% \]

% \subsubsection{Proportionate Relationship}
% The AULTC in units $V$ ($\text{AULTC}^{V}$) is approximately proportional to the AULTC in units $U$ ($\text{AULTC}^{U}$) via the ratio of their normalization factors:
% \[
% \text{AULTC}^{V} \approx \left( \frac{\log(S_{\text{max}}^{U})}{\log(c) + \log(S_{\text{max}}^{U})} \right) \times \text{AULTC}^{U}
% \]

% \subsubsection{Derivation}
% For this derivation, we assume that approximation $\log(1 + S_{\text{max}}^{U}) \approx \log(S_{\text{max}}^{U})$ holds.
% For units $U$, we have the approximation of the normalizing factor $\gamma_U$:
% \[
% \gamma_U = \log(1 + S_{\text{max}}^{U}) \approx \log(S_{\text{max}}^{U})
% \]
% $(\text{for large } S_{\text{max}}^{U})$. For units $V$: 
% \[
% \gamma_V = \log(1 + c S_{\text{max}}^{U}) \approx \log(c) + \log(S_{\text{max}}^{U})
% \]
% $\text{for large } c S_{\text{max}}^{U}$.
% Then, the ratio of the normalization factors is:
% \[
% \frac{\log(S_{\text{max}}^{U})}{\log(c) + \log(S_{\text{max}}^{U})}
% \]
% This ratio scales the AULTC when switching units.

% With these approximations, the proportionate relationship is:
% \[
% \text{AULTC}^{V} \approx \left( \frac{\log(S_{\text{max}}^{U})}{\log(c) + \log(S_{\text{max}}^{U})} \right) \times \text{AULTC}^{U},
% \]
% which allows conversion of AULTC between any units $U$ and $V$ with a conversion factor $c$, provided $S_{\text{max}}^{U}$ is sufficiently large.

% \subsubsection{Example}

% Here is an example proportionality of the AULTC when converting from units of hours to units of seconds when the $S_\text{max}$ is 1 year.

% For $ S_{\text{max}}^{\text{hr}} = 8760 $ hours (1 year):
% \begin{align*}
% &\log(8760) \approx 9.07, \\
% &\log(3600 \times 8760) = \log(31,536,000) \approx 17.27, \\
% &\text{Ratio:}
%    \frac{9.07}{17.27} \approx 0.525
% \end{align*}
% Thus:
% \[
% \text{AULTC}^{\text{sec}} \approx 0.525 \times \text{AULTC}^{\text{hr}}
% \]